%% ---------------------------------------------------------------------
\section{扭度 3 与单自旋不对称}\label{sec:power-twist3}

横向单自旋不对称 $A_N=[\dd\sigma(S_{\mathrm T})-\dd\sigma(-S_{\mathrm T})]/[\dd\sigma(S_{\mathrm T})+\dd\sigma(-S_{\mathrm T})]$ 在领头扭度的共线因子化中为零：它正比于“朴素 T 奇”的结构 $\mathbf S_{\mathrm T}\cdot(\mathbf P_A\times\mathbf P_h)$，既要求振幅之间有相对相位，又要求螺旋度翻转，而无质量夸克的手征在微扰论中守恒；微扰 QCD 对产生夸克横向极化的预言在标度极限下为零\cite{Kane:1978nd}。当 $P_{h\mathrm T}\ll Q$ 时，相应的机制是 TMD 中的 Sivers 函数（第~14~章）；当 $P_{h\mathrm T}\sim Q$ 时，不对称是扭度 3 效应：强子经夸克—胶子—夸克关联与硬部分相连，相位来自硬部分中传播子的极点。基本对象是 Efremov--Teryaev--Qiu--Sterman（ETQS）函数 $T_F(x_1,x_2)$（\cref{def:tmd-etqs}）\cite{Efremov:1981sh,EfremovTeryaev1985,QiuSterman1991c}，其对角值 $T_F(x,x)$ 称为胶子极点矩阵元；在两种描述都适用的区域 $\Lambda\ll q_{\mathrm T}\ll Q$，二者在单圈上一致\cite{Ji:2006ub}（\cref{prop:tmd-twist3-matching}）。

\begin{proposition}[SIDIS 型过程中的扭度 3 共线因子化 \tagP]\label{prop:power-twist3-sidis}
在至多一个初态强子的过程中（如 $P_{h\mathrm T}\sim Q$ 的半单举 DIS），或两个初态强子而被测末态为色单态、其余末态单举求和的过程中（如 $q_{\mathrm T}\sim Q$ 的 Drell--Yan）——统称“SIDIS 型”——横向单自旋不对称在扭度 3 上具有共线因子化形式：自旋相关截面写成硬函数与 $T_F(x,x)$（及其导数）、其他扭度 3 关联函数和未极化分布或碎裂函数的卷积，余项为更高幂次。具体地：(a) 这一形式在特定过程中已在单圈与 NLO 上验证\cite{Ji:2006ub,Vogelsang:2009pj}；(b) $T_F(x,x)$ 的标度演化已经确定：最初的结果互有出入（\cite{Braun:2009mi}与\cite{Kang:2008ey}），分歧的来源后来被找出，各组结果得到调和\cite{Kang:2012em,Schafer:2012ra}；(c) $T_F(x,x)$ 与 SIDIS 中 Sivers 函数的第一个横动量矩由
\begin{equation}\label{eq:power-TF-sivers}
g\,T_F(x,x)=\mp\int\!\dd^2k_{\mathrm T}\,\frac{\mathbf k_{\mathrm T}^2}{M}\,f^{\perp q}_{1T}(x,k_{\mathrm T})\Big|_{\mathrm{SIDIS}}
\end{equation}
相联系，符号与归一化依赖于约定\cite{Boer:2003cm}（本书的约定见\cref{prop:tmd-twist3-matching}）。
非微扰地位：\cref{asm:intro-bridge}（MA1、MA2）。
\uses{def:tmd-etqs,prop:tmd-twist3-matching,asm:intro-bridge}
\end{proposition}

\begin{proofsketch}
1.~把硬部分对部分子横动量展开到一阶，并加上多一条共线胶子接入硬部分的图，用 Ward 恒等式组合成规范不变的夸克—胶子—夸克关联（与\cref{prop:power-dis-twist4-x}同一方法，停在一阶）；所需相位来自硬部分中传播子的极点，把关联函数钉在 $T_F(x,x)$ 等点上（\cref{ex:power-sgp}）。2.~单圈匹配与 NLO 计算见\cite{Ji:2006ub,Vogelsang:2009pj}，演化见\cite{Braun:2009mi,Kang:2008ey,Kang:2012em,Schafer:2012ra}。3.~\eqref{eq:power-TF-sivers}在算符层面由 Sivers 函数的 $k_{\mathrm T}$ 矩导出\cite{Boer:2003cm}；在 $D_\mu=\partial_\mu+\ii gA_\mu$、$\epsilon^{0123}=+1$、Trento 约定下右边取负号\cite[式~(10)]{Kang:2011hk}，换到本书的 $D_\mu=\partial_\mu-\ii gA_\mu$ 而保持 $T_F$ 的定义式不变时 $gG^{\mu\nu}$ 变号，右边随之反号。完整论证见所引文献；不存在全阶论证。
\end{proofsketch}

\paragraph{地位与缺口}
证据都是有限阶、过程特定的（单圈与 NLO 计算、单圈匹配、领头阶演化），没有次领头幂次的区域分析或积分范数估计，故为 \tagP。\eqref{eq:power-TF-sivers}受紫外减除方案的影响，不是唯一的恒等式\cite[第~II~节]{Kang:2011hk}。唯象上，SIDIS 中拟合的 Sivers 矩与 $pp$ 中拟合的 $T_F(x,x)$ 大小相当而符号相反（“符号失配”）\cite{Kang:2011hk}；同时拟合 SIDIS、Drell--Yan、$e^+e^-$ 与 $pp$ 数据的全局分析认为各种单自旋不对称有共同起源，并认为这指向失配的解决\cite{Cammarota:2020qcw}。升级途径：扭度 3 区域分析的全阶论证，以及胶子极点矩阵元的方案无关定义。

\begin{conjecture}[强子—强子碰撞中单举强子单自旋不对称的扭度 3 因子化 \tagC]\label{conj:power-twist3-pp}
$p^\uparrow p\to hX$ 中大横动量单举强子\cite{QiuSterman1991c,Qiu:1998ia}与 $p^\uparrow p\to\gamma X$ 中直接光子\cite{QiuSterman1992}的单自旋不对称，在扭度 3 上具有\cref{prop:power-twist3-sidis}那样的共线因子化形式。
非微扰地位：\cref{asm:intro-bridge}（MA1、MA2）。
\uses{prop:power-twist3-sidis,conj:dy-coloured-final,asm:intro-bridge}
\end{conjecture}

\paragraph{地位与缺口}
末态含大横动量的带色部分子（直接光子的领头阶子过程 $qg\to\gamma q$、$q\bar q\to\gamma g$ 中反冲的部分子），旁观者与它之间的 Glauber 交换未被全阶分析；即使在领头扭度，带色末态的共线因子化也没有全阶论证（\cref{conj:dy-coloured-final}，\tagC），扭度 3 的陈述继承这一地位。“符号失配”\cite{Kang:2011hk}及其全局分析\cite{Cammarota:2020qcw}检验的是唯象框架，不是因子化本身。升级途径：首先需要领头扭度带色末态因子化的全阶论证。

\begin{example}[最低阶的软胶子极点机制]\label{ex:power-sgp}
考虑\cref{fig:twist3-ssa}中 $A^\uparrow B\to\gamma X$ 的 Compton 子过程，$P_A=(P_A^+,0^-,\mathbf0_T)$。$T_F$ 提供的胶子（动量 $x_gP_A$，$x_g=x_2-x_1$）在振幅一侧接到末态夸克 $c$ 上，中间夸克传播子携带 $p_c-x_gP_A$；由 $p_c^2=P_A^2=0$，
\begin{equation}\label{eq:power-sgp-den}
(p_c-x_gP_A)^2+\ii0=-2\,p_c\cdot P_A\,(x_g-\ii0),\qquad \frac{1}{x_g-\ii0}=\mathrm P\frac{1}{x_g}+\ii\pi\,\delta(x_g),
\end{equation}
其中 $p_c\cdot P_A>0$。$\delta$ 函数把胶子钉在 $x_g=0$（“软胶子极点”），给出 $\ii\pi$ 与 $T_F(x,x)$；这个 $\ii$ 与 $T_F$ 中的 $\epsilon^{S_{\mathrm T}\sigma n_An_B}$ 相乘，使自旋相关截面为实且非零——树图上传播子的 $\pm\ii0$ 是唯一的相位来源。若胶子接到来自 $B$ 的入射部分子上（初态相互作用），分母为 $2p_b\cdot P_A(x_g+\ii0)$。一般地，在程函极限下（$k$ 为流入夸克线的胶子动量），出射与入射夸克分别给出
\begin{equation}\label{eq:power-sgp-eik}
\bar u(p)\,\ii g\gamma^\mu t^a\frac{\ii(\sla p-\sla k)}{(p-k)^2+\ii0}\to\bar u(p)\frac{g\,t^ap^\mu}{p\cdot k-\ii0},\qquad
\frac{\ii(\sla p+\sla k)}{(p+k)^2+\ii0}\,\ii g\gamma^\mu t^au(p)\to-\frac{g\,t^ap^\mu}{p\cdot k+\ii0}u(p)
\end{equation}
（与\cref{prop:prelim-wilson-rules}一致）。相对于主值部分，极点部分在末态相互作用中是 $+\ii\pi\delta$，在初态相互作用中是 $-\ii\pi\delta$；前者构成未来指向的 Wilson 线（$1/(k^+-\ii0)$），后者构成过去指向的线（$1/(k^++\ii0)$），二者之差为 $2\pi\ii\,\delta(k^+)$。因此单自旋不对称在半单举 DIS（末态相互作用）与 Drell--Yan（初态相互作用）中符号相反，这正是\cref{thm:tmd-sivers-sign}（\tagR，$R_{\mathrm{def}}$）的符号反转在扭度 3 语言中的对应。在 $pp\to hX$ 中，多种接法带有不同的色因子与 $\pm\ii0$，过程依赖性更复杂，这是\cref{conj:power-twist3-pp}停留在猜想层次的原因之一。（上述代数均用符号与数值计算核对过。）
\uses{prop:power-twist3-sidis,thm:tmd-sivers-sign}
\end{example}

%%FIG:fig_twist3_ssa%%
